%!TEX root = ../deepsmooth.tex
 
\section{Methodology} \label{sec:smoo}

\subsection{Model and loss function} \label{sec:archlearn}

\textbf{Explanatory and target variables.}
At a given time we observe triplets $(\sigma_i, \, k_i, \, \tau_i)\in\datamkt$
where $\sigma_i$ is the market implied volatility (the target/response), and $(k_i, \tau_i)$ are the log moneyness and the time to maturity (the features/explanatory variables).
In addition, we complement the sample with synthetic pairs $(k_i, \, \tau_i)\in \datacalbut \cup\dataasy$ used to control the arbitrage opportunities and the model asymptotic behavior (see~\Cref{sec:noarb}).

\textbf{Implied volatility model.}
Our model for the total variance and implied volatility is given by
\begin{align}\label{eq:iv_ann}
\wth(k,\tau) &= \wnn\big(k,\tau;\theta_1\big) \times \wpr\big(k,\tau;\theta_2\big)  \quad \mbox{and} \quad
\ivth(k,\tau) = \sqrt{\wth(k,\tau)/\tau}
\end{align} 
for the parameters $\theta=\{\theta_1, \theta_2\}$, where
$\wnn$ and $\wpr$ are the NN and prior models described below.

\textbf{Prior model.}
$\wpr:\R^2\mapsto\R$ is a prior model with parameters $\theta_2$.
An implied volatility model without prior is obtained by setting $\wpr\equiv 1$ so that $\wth \equiv \wnn $. 
The prior model choice is useful to ensure that the model generalization is compliant with a prescribed preferred behavior.
Essentially, the BS prior is a parameter-free model matching the implied volatility's at-the-money (ATM) term-structure.
As for the SSVI prior, it improves on the BS model by capturing both the smile and the skew of the surface.
Both models are described in Appendix~\ref{sec:bsprior} and~\ref{sec:ssviprior} respectively.

\textbf{NN model.}
$\wnn:\R^2\mapsto\R$ is a standard feedforward multilayer neural network, namely
\begin{align}
\wnn(k, \tau;\theta_1) &=  \comp_{i=1}^{n+1} f_i^{W_i, b_i} (k, \tau) \mbox{ with } 
f_i(x) = \begin{cases} g_i(W_i \, x + b_i) & i< n + 1
\\ \alpha \left(1 + \tanh \left(W_{n+1} \, x + b_{n+1} \right)\right) & i=n+1 \end{cases} \label{eq:ffnn_def2}
\end{align}
with $g_i$ an activation function, $\theta_1=\{W_1,\,b_1,\,W_2,\,b_2,\dots, \alpha\}$ the set of weight matrices and bias vectors, $\alpha$ a scaling parameter letting $\wnn$ take values in $[0, \alpha]$, and $n$ the number of hidden layers.
The functional choice $1+\tanh$ in the last layer is not critical, yet it appeared to perform slightly better than the sigmoid function on a limited number of trials.
Note that one can initialize the NN to take output value one, that is no NN correction, by setting $\alpha=1$ and $W_{n+1}=b_{b+1}=0$.

\textbf{Loss function.}
We fit the network parameters and prior parameters $\theta$ by minimizing the loss function
\begin{equation}\label{eq:loss}
\Lcal (\theta) = \Lcal_0(\theta) + \sum_{j=1}^6 \; \lambda_j \, \Lcal_{{\rm C} j}(\theta)
\end{equation}
where the term $\Lcal_0(\theta)$ is a prediction error cost, the terms $\Lcal_{{\rm C} j}(\theta)$ for $j=1,\dots,6$ materialize soft constraints aiming to ensure that the shape of $\{\wth(k, \tau);\, (k, \tau) \in\R\times\R_+\}$ is indeed a sensible implied volatility surface, and $\lambda_i$ for $j=1,\dots,6$ are the corresponding penalty weights.
Note that some parameters of the prior model may also be calibrated.

We let the prediction error be the sum of the root-mean-squared-error (RMSE) and the mean-absolute-percentage-error (MAPE),
\[
\Lcal_0(\theta) = \sqrt{1/|\datamkt|\sum_{(\sigma_i,k_i, \tau_i)\in\datamkt} (\sigma_i - \ivth(k_i, \tau_i))^2} + \;1/|\datamkt| \sum_{(\sigma_i,k_i, \tau_i)\in\datamkt}  \vert \sigma_i - \ivth(k_i, \tau_i) \vert/\sigma_i,
\]
so as to penalize both absolute and relative errors.
The IV values far out of the money can take values significantly larger than at the money.
We found that the above loss function results in general a balanced allocation of prediction errors across the different moneynesses.

\begin{remark}[Soft versus hard constraints]
An alternative approach to impose shape constraints on the mapping $\wth$ is to hard-wire them into the neural network architecture, as in~\cite{dugas2001incorporating} for example.
However, hard constraints are difficult to impose on multilayer neural networks, may reduce the neural network's flexibility, and may lead to more challenging leaning routines, see~\cite{marquez2017imposing}.
\end{remark}

\subsection{No-arbitrage conditions and synthetic grid} \label{sec:noarb}

We explain how each of the constraints/conditions in~\Cref{prop:noAO} can be handled either by refining the architecture of the neural network, or by adding a penalty term to the loss function~\eqref{eq:loss}.
Note that the conditions \textbf{\ref{cons:posi}--\ref{cons:vmat}} are in principle satisfied by design of the NN.
The mapping $\wth$ is twice differentiable as long as the activation functions $g_i$ and $g_{n+1}$ (as well as the prior model) are twice differentiable, in which case\textbf{ \ref{cons:smoo}} is satisfied.
An example of valid activation function is the SoftPlus given by $\ln(1+\exp(x))$\footnote{We conjecture that one could also use ReLU activation functions with adjusted constraint conditions.}.
Hence we set $\Lcal_{\rm C1} \equiv \Lcal_{\rm C2}\equiv\Lcal_{\rm C3}\equiv 0$.

As for the other three constraints, we control for
\begin{itemize}
\item calendar arbitrages with $\Lcal_{\rm C4} (\theta) = 1/|\datacalbut| \sum_{(k_i,\tau_i)\in\datacalbut} \max\left(0, \, - \ell_{\rm cal}(k_i,\,\tau_i) \right),$
\item butterfly arbitrages with $\Lcal_{\rm C5} (\theta) = 1/|\datacalbut| \sum_{(k_i,\tau_i)\in\datacalbut}  \max\left(0, \, -\ell_{\rm but}(k_i,\,\tau_i)\right)$,
\item and the asymptotic behavior with $\Lcal_{\rm C6} (\theta) = 1/|\dataasy| \sum_{(k_i,\tau_i)\in\dataasy} \left\lvert \partial^2 \wth(k_i,\tau_i)/\partial k \partial k \right\rvert$,
\end{itemize}
where 
\begin{align*}
\datacalbut &= \big\{(k,\tau) \, : \, k\in \big\{ x^3 \,:\, x \in \big\lbrack -(-2 k_{\min})^{1/3}, \, (2k_{\max})^{1/3} \big\rbrack_{100} \big\}, \,\tau \in \Tcal\big\} \\
\dataasy &= \big\{(k,\tau) \, : \, k\in \big\{6k_{\min}, \, 4k_{\min} , \, 4k_{\max}, \, 6k_{\max} \big\}, \,\tau \in \Tcal\big\}\\
\Tcal &= \big\{ \exp(x) \,:\, x \in \big\lbrack \log(1/365), \, \max( \log(\tau_{\max}+1)) \big\rbrack_{100} \big\}
\end{align*}
 $k_{\min} = \min(\datamkt^k)$, $k_{\max} = \max(\datamkt^k)$, $\tau_{\max} = \max(\datamkt^\tau)$, $\lbrack a,b \rbrack_{x}$ indicates an equidistant set of $x$ points between $a$ and $b$, and $\datamkt^\tau$ and $\datamkt^k$ are the sets of unique time to maturity and forward log moneyness in $\datamkt$.
Note that it should always be that $\min(\datamkt^k)<0$ and $\max(\datamkt^k)>0$.
The motivation for the above transformations is to obtain a denser grid around the money and for short maturities.
The particular parametric choices for the above sets do not appear critical as long as they are sufficiently dense and cover the regions of interest.

%As for the other three constraints, we use $\Lcal_{\rm C4} (\theta) = 1/|\datacal| \sum_{(k_i,\tau_i)\in\datacal} \max\left(0, \, - \ell_{\rm cal}(k_i,\,\tau_i) \right)$ to prevent
% calendar arbitrage,
%$\Lcal_{\rm C5} (\theta) = 1/|\databut| \sum_{(k_i,\tau_i)\in\databut}  \max\left(0, \, -\ell_{\rm but}(k_i,\,\tau_i)\right)$ for butterfly arbitrage,
%and $\Lcal_{\rm C6} (\theta) = 1/|\dataasy| \sum_{(k_i,\tau_i)\in\dataasy} \left\lvert \partial^2 \wth(k_i,\tau_i)/\partial k \partial k \right\rvert$ for the large moneyness behavior.


Note that the loss $\Lcal_{\rm C6}$ guarantees that $\omega$ is asymptotically linear and should in principle be refined to control the coefficient value to be less than 2, as discussed in~\cite{lee2004moment}.
However $\Lcal_{\rm C6}$ has better training performance and proved to be compliant with~\ref{cons:asym} in our post training verifications.
Furthermore, we only work with finite moneyness in practice and imposing an asymptotic condition on large but finite moneyness values may unnecessarily constraint the model.

\subsection{Model training}

The training procedure and the parameters choice are described in details in Appendix~\ref{sec:train}.